\scanpage{11}
\uline{Proposition} a) \inserted{$G$ \deleted{\unsure{disons}} \unsure{groupe} algébrique} Il existe un \deleted{\unsure{hom|hom.}}\inserted{hom} [\uline{épimorphisme}]
\[
G\longrightarrowH
\]
$H$ de t. m., tel que \deleted{\unsure{tt|$H$} \unsure{autre}} hom de $G$ dans un gpe de t. m. se factorise de façon unique par $H$.
% in the left margin
\diagram{\begin{tikzcd}[ampersand replacement=\&,row sep=0.6ex,column sep=0.5em]
 R \& \subset \& Z \& \deleted{\unsure{G}}\arrow[dd] \\
 \cap \& \& \cap \& {} \\
 T \& \subset \& G\arrow[d] \& {}\\
 {} \& {} \& \Gamma \& {}
\end{tikzcd}}

b) Lorsque \uline{$k$ est parfait}, la formation de $H$ est compatible aux \deleted{\unsure{tt|$H$}} extensions de la base.
\margin{$G$}

c) Lorsque $G$ est connexe \inserted{et lisse}, $H$ est un tore.
% in the left margin beside N.B.
\diagram{\begin{tikzcd}[ampersand replacement=\&,row sep=1em,column sep=1.7em]
 G\arrow[r,"\text{épim}"]\arrow[d,dash]\arrow[dr,phantom,"\deleted{\unsure{a}}" description] \& H\\
 G^{\circ}_{\text{réd}}\arrow[r,dash] \& H^{\circ}\arrow[u]
\end{tikzcd}}
\[
G/G^{\circ}_{\text{réd}}\unsure{\mathord{-}|\to}H/H^{\circ}
\]
\uline{N.B.} Si $k$ est imparfait, $\unsure{=|\simeq}$ si $G$ est commutatif, \inserted{\unsure{lisse sur le|linéaire et}\ \unsure{même|corps},} la formation de $H$ n'est en général pas compatible avec

\deleted{\uline{Corollaire}}\inserted{Proposition} $G$ un groupe\deleted{\unsure{$k$}} connexe \deleted{lisse} sur $k$, $R$ un \deleted{gpe de \unsure{t. b.|t. m.}}\inserted{tore}. Alors tte extension de $G$ par un $R$ qui splitte en tant que fibré principal homogène \uline{splitte}.
\[
\deleted{H^{2}_{\text{Hochschild}}(G,R)\simeq}
\]
\[
\left.\begin{aligned}
0\to C^{\bullet}(e,R)\to C^{\bullet}(H,R)&\longrightarrow
\deleted{E}\operatorname{Hom}_{\mathrm{gr}}(H^{\bullet},R)\mathrel{-}0\\
&H^i(\operatorname{Hom}_{\mathrm{gr}}(H^{\bullet},R))
\end{aligned}\right]
\]
\margin{En effet \inserted{\unsure{l'ext. est centrale|d'ext. centrale}} \unsure{Conjuguant|Composant} les deux résultats qui précèdent, elle provient d'une extension de $H$ par $R$, donc splitte.}
% archive sheet number: 10
